%! Mode:: "TeX:UTF-8"
%! TEX program = xelatex
\PassOptionsToPackage{quiet}{xeCJK}
\PassOptionsToPackage{table,xcdraw}{xcolor}
\documentclass[withoutpreface,bwprint]{cumcmthesis}
\usepackage{tcolorbox}
\usepackage{newtxmath}  % Times 风格数学字体
%\usepackage{newtxmath}  %使用Times New Roman 字体的方法
%\usepackage{mathptmx}%使用Times New Roman 字体的方法
%\usepackage{txfonts}%使用Times New Roman 字体的方法,全部注释掉也行。
\setlength{\heavyrulewidth}{1.5pt} % 设置顶部和底部线条的粗细
\setlength{\lightrulewidth}{1pt} % 设置中间线条的粗细
\usepackage{threeparttable}
\usepackage[numbers,sort&compress]{natbib}  % 文献管理宏包
\definecolor{lightblue1}{RGB}{173,216,230} % 浅蓝（示例色）
\definecolor{lightblue2}{RGB}{119,191,219} % 中蓝（示例色）
\definecolor{lightblue3}{RGB}{70,130,180}  % 深蓝（示例色）

\newcolumntype{C}{>{\centering\arraybackslash}X}
\newcolumntype{R}{>{\raggedleft\arraybackslash}X}
\newcolumntype{L}{>{\raggedright\arraybackslash}X}

% 定义颜色
\definecolor{codegreen}{rgb}{0,0.5,0} % 更深的绿色
\definecolor{codegray}{rgb}{0.5,0.5,0.5}
\definecolor{codepurple}{rgb}{0.58,0,0.82}
\definecolor{backcolour}{rgb}{1,1,1} % 白色背景
\definecolor{titlebg}{rgb}{0.9,0.9,0.9} % 标题浅灰色背景
\definecolor{blackframe}{rgb}{0,0,0} % 深黑色边框

% 定义基本代码样式
\lstdefinestyle{mystyle}{
	backgroundcolor=\color{backcolour},   % 白色背景
	commentstyle=\color{codegreen},       % 注释使用绿色
	keywordstyle=\color{blue},            % 关键字使用蓝色
	numberstyle=\tiny\color{codegray},
	stringstyle=\color{red},              % 字符串使用红色
	basicstyle=\ttfamily\footnotesize,
	breakatwhitespace=false,         
	breaklines=true,                 
	captionpos=b,                    
	keepspaces=true,                 
	numbers=left,                    
	numbersep=5pt,                  
	showspaces=false,                
	showstringspaces=false,
	showtabs=false,                  
	tabsize=4,
	frame=single,                    % 单线边框
	framesep=0pt,                    % 边框与内容无间距
	framerule=1pt,                   % 边框线宽
	rulecolor=\color{blackframe},    % 深黑色边框
	xleftmargin=0pt,                 % 左边框无额外间距
	xrightmargin=0pt,                % 右边框无额外间距
	morestring=[b]{"}                % 识别双引号字符串
}



% 为 Python 设置特定样式
\lstdefinestyle{Pythonstyle}{
	style=mystyle,
	language=Python,
	morecomment=[l]{\#},             % 识别 # 开头的行注释
	morekeywords={import, as, def, class, for, while, if, elif, else, 
		return, try, except, finally, with, in, is, not, and, or,
		print, lambda, None, True, False, pass, break, continue,
		pd, np, cp, random, math, Prophet, datetime, timedelta,
		plt, cross_validation, performance_metrics, LinearRegression,
		mean_squared_error, mean_absolute_error, sm, read_excel,
		groupby, merge, loc, reset_index, pivot_table, to_excel}
}

% 自定义标题框样式
\newenvironment{codeheader}[1]{%
	\begin{tcolorbox}[
		colback=titlebg,
		colframe=blackframe,
		sharp corners,
		boxrule=1pt,
		left=0pt,
		right=0pt,
		top=0pt,
		bottom=0pt,
		boxsep=5pt,
		before skip=10pt,
		after skip=10pt
		]%
		\textbf{#1}
	\end{tcolorbox}
	\vspace{-5pt}
	\hrule height 1pt\relax
	\vspace{5pt}
}{}

\title{}  % 论文标题
%%%%%%%%%%%%%%%%%%%%%%%%%%%%%%%%%%%%%%%%%%%%%%%%%%%%%%%%
%% 正文

\begin{document}
	
	\maketitle
	\begin{abstract}
		
		\textbf{关键词：{}}
	\end{abstract}
	%%%%%%%%%%%%%%%%%%%%%%%%%%%%%%%%%%%%%%%%%%%%%%%%%%%%%%%%%%%%% 
	
	% \tableofcontents  % 目录
	% \newpage
	
	%%%%%%%%%%%%%%%%%%%%%%%%%%%%%%%%%%%%%%%%%%%%%%%%%%%%%%%%%%%%%  
	%%%%%%%%%%%%%%%%%%%%%%%%%%%%%%%%%%%%%%%%%%%%%%%%%%%%%%%%%%%%%
	%% 参考文献
	\section{问题重述}
	\subsection{问题背景}
	\subsection{问题提出}
	\section{问题分析}
	\subsection{问题一分析}
	\subsection{问题二分析}
	\subsection{问题三分析}
	\subsection{问题四分析}
	\section{模型假设}
	\section{符号说明}
	\section{数据预处理}
	\section{问题一模型的建立与求解}
	\section{问题二模型的建立与求解}
	\section{问题三模型的建立与求解}
	\section{问题四模型的建立与求解}
	\section{模型灵敏度分析}
	\section{模型评价与推广}
	
	\newpage
	\begin{thebibliography}{99}
		\bibitem{a}
	\end{thebibliography}
	
	\newpage
	%%%%%%%%%%%%%%%%%%%%%%%%%%%%%%%%%%%%%%%%%%%%%%%%%%%%%%%%%%%%%
	%% 附录
	\begin{appendices}
		\section{支撑材料文件列表}
		\renewcommand{\baselinestretch}{1}%和下面那个并不直接冲突，改不改都行
		\begin{table}[H]
			\tabcolsep=0.7cm
			\renewcommand\arraystretch{1.5}  % 正确的应该是 1.5 而不是 -1.5，以避免行间距被压缩，和上面不冲突，如果有就是在它的基础上。
			\centering
			\vspace{1ex}
			\begin{tabular}{ll}
				\toprule[1.5pt]
				\textbf{文件名} & \textbf{功能描述} \\
				\midrule[1pt]
				
				\bottomrule[1.5pt]
			\end{tabular}
		\end{table}
		
		
		
		
		\section{代码}%%%captionpos=t让名称置顶出现
		\begin{codeheader}{附录1 数据预处理代码}
		\end{codeheader}
		\begin{lstlisting}[style=Pythonstyle]
			
		\end{lstlisting}
		
		
		\begin{codeheader}{附录2 问题一的代码}
		\end{codeheader}
		\begin{lstlisting}[style=Pythonstyle]
			
		\end{lstlisting}
		
		
		\begin{codeheader}{附录3 问题二代码}
			
		\end{codeheader}
		\begin{lstlisting}[style=Pythonstyle]
			
		\end{lstlisting}
		
		\begin{codeheader}{附录4 问题三代码}
		\end{codeheader}
		\begin{lstlisting}[style=Pythonstyle]
			
		\end{lstlisting}
		
		
		
		
	\end{appendices}
	
\end{document}